\documentclass{article}
\usepackage{amsmath,amssymb,amsthm}
\usepackage{algorithm}
\usepackage{algpseudocode}
\usepackage{hyperref}
\newtheorem{theorem}{Theorem}
\newtheorem{lemma}{Lemma}
\newtheorem{assumption}{Assumption}
\newcommand{\prox}{\operatorname{prox}}
\newcommand{\R}{\mathbb{R}}

\begin{document}

\section*{Algorithm Name}
\textbf{Proximal Gradient Method with Gradient Momentum (PGMGM)}

\section*{Mathematical Formulation}
Let 
\[
F(x)=f(x)+g(x),\qquad x\in\R^{d},
\]
where  

\begin{itemize}
    \item $f:\R^{d}\!\to\!\R$ is convex and $L$–smooth, i.e. 
    \[
    \|\nabla f(u)-\nabla f(v)\|\le L\|u-v\|,\qquad\forall u,v\in\R^{d}.
    \]
    \item $g:\R^{d}\!\to\!\R\cup\{+\infty\}$ is proper, closed and convex.
\end{itemize}

For a stepsize $\eta>0$ and a momentum parameter $\beta\in[0,1)$ we define the \emph{gradient momentum} at iteration $k\ge0$ by  
\[
v_{k}:=\nabla f(x_{k})+\beta\bigl(\nabla f(x_{k})-\nabla f(x_{k-1})\bigr),
\]
with the convention $\nabla f(x_{-1})\equiv 0$.  
The update rule of PGMGM is
\[
\boxed{\;
x_{k+1}= \prox_{\eta g}\bigl(x_{k}-\eta v_{k}\bigr),\qquad k=0,1,2,\dots
\;}
\tag{1}
\]
The proximal operator is
\[
\prox_{\eta g}(z):=\arg\min_{u\in\R^{d}}\Bigl\{ g(u)+\frac{1}{2\eta}\|u-z\|^{2}\Bigr\}.
\]

\section*{Pseudocode}
\begin{algorithm}[H]
\caption{PGMGM}
\begin{algorithmic}[1]
\Require stepsize $\eta\in(0,1/L]$, momentum $\beta\in[0,1)$, initial point $x_{0}\in\R^{d}$.
\State Set $x_{-1}\gets x_{0}$.
\For{$k=0,1,2,\dots$}
    \State Compute gradients $g_{k}:=\nabla f(x_{k})$, $g_{k-1}:=\nabla f(x_{k-1})$.
    \State Form momentum‑augmented gradient 
    \[
    v_{k}\gets g_{k}+\beta\,(g_{k}-g_{k-1}) .
    \]
    \State Proximal step
    \[
    x_{k+1}\gets \prox_{\eta g}\bigl(x_{k}-\eta v_{k}\bigr).
    \]
\EndFor
\end{algorithmic}
\end{algorithm}

\section*{Dry Run Example}
Consider the scalar problem
\[
f(w)=(w-3)^{2},\qquad g\equiv0,
\]
so $F(w)=f(w)$.  
We take $\eta=0.1$, $\beta=0.5$, $w_{0}=0$ and set $w_{-1}=w_{0}=0$.

\begin{center}
\begin{tabular}{c|c|c|c|c}
$k$ & $w_{k}$ & $\nabla f(w_{k})=2(w_{k}-3)$ & $v_{k}$ & $w_{k+1}=w_{k}-\eta v_{k}$\\ \hline
0 & $0$ & $-6$ & $v_{0}= -6+\beta(-6-0)= -6-3 = -9$ & $w_{1}=0-0.1(-9)=0.9$\\[4pt]
1 & $0.9$ & $2(0.9-3)=-4.2$ & $v_{1}= -4.2+0.5(-4.2+6)= -4.2+0.9 = -3.3$ & $w_{2}=0.9-0.1(-3.3)=1.23$\\[4pt]
2 & $1.23$ & $2(1.23-3)=-3.54$ & $v_{2}= -3.54+0.5(-3.54+4.2)= -3.54+0.33 = -3.21$ & $w_{3}=1.23-0.1(-3.21)=1.551$\\
\end{tabular}
\end{center}
The objective values are $f(w_{0})=9$, $f(w_{1})\approx5.61$, $f(w_{2})\approx4.07$, $f(w_{3})\approx2.96$, showing monotone decrease.

\section*{Assumptions}
\begin{assumption}[Convexity and Smoothness of $f$]\label{as:convexsmooth}
$f$ is convex and $L$‑smooth with constant $L>0$:
\[
f(u)\le f(v)+\langle\nabla f(v),u-v\rangle+\frac{L}{2}\|u-v\|^{2},\qquad\forall u,v\in\R^{d}.
\]
\end{assumption}
\begin{assumption}[Convexity of $g$]\label{as:gconvex}
$g$ is proper, closed and convex; its proximal operator $\prox_{\eta g}$ is well defined for any $\eta>0$.
\end{assumption}
\begin{assumption}[Stepsize and Momentum]\label{as:params}
The stepsize satisfies $0<\eta\le 1/L$ and the momentum parameter obeys
\[
0\le\beta<\frac{1-\eta L}{\eta L}\quad\text{(in particular, }\beta<1\text{)}.
\]
\end{assumption}
\begin{assumption}[Existence of a Minimizer]\label{as:opt}
There exists at least one $x^{*}\in\arg\min_{x}F(x)$ and $F(x^{*})$ is finite.
\end{assumption}

\section*{Main Convergence Theorem}
\begin{theorem}[Sublinear $O(1/k)$ Convergence]\label{thm:sublinear}
Under Assumptions \ref{as:convexsmooth}–\ref{as:opt}, let $\{x_{k}\}$ be generated by \eqref{1}.  
Define the constant
\[
C:=\frac{1}{2\eta}\|x_{0}-x^{*}\|^{2}.
\]
Then for all $k\ge1$,
\[
F(x_{k})-F(x^{*})\;\le\;\frac{C}{k}.
\tag{2}
\]
Consequently $\displaystyle\lim_{k\to\infty}F(x_{k})=F(x^{*})$.
\end{theorem}

\section*{Theoretical Properties}
\begin{itemize}
    \item \textbf{Stability.} The bound \eqref{2} holds for any $\eta\le1/L$; decreasing $\eta$ reduces the constant $C$ but slows down progress per iteration.
    \item \textbf{Momentum Sensitivity.} The admissible range for $\beta$ shrinks as $\eta$ approaches $1/L$; when $\beta=0$ the method reduces to the classical proximal gradient method with its standard $O(1/k)$ rate.
    \item \textbf{Comparison.} Compared with vanilla proximal gradient, PGMGM enjoys the same asymptotic rate while often exhibiting a smaller constant $C$ in practice because the momentum term reduces the distance $\|x_{k}-x^{*}\|$ faster.
\end{itemize}

\section*{Discussion}
The result shows that adding a simple gradient momentum does not deteriorate the classical $O(1/k)$ guarantee for composite convex optimisation. The proof follows the standard proximal‑gradient analysis with an extra term that is controlled by the stepsize–momentum condition \eqref{as:params}. The theorem can be extended to stochastic settings by incorporating unbiased gradient estimators and variance bounds, at the price of an additional $\mathcal{O}(1/\sqrt{T})$ term.

\section*{Preliminaries (Definitions and Lemmas)}
\begin{lemma}[Three‑point identity for the proximal operator]\label{lem:prox}
For any $z\in\R^{d}$ and any $u\in\R^{d}$,
\[
\langle \prox_{\eta g}(z)-z,\,u-\prox_{\eta g}(z)\rangle
    +\eta\bigl(g(\prox_{\eta g}(z))-g(u)\bigr)\le 0.
\]
\end{lemma}
\begin{proof}
The optimality condition of the proximal definition yields
\[
0\in \partial g(\prox_{\eta g}(z))+\frac{1}{\eta}\bigl(\prox_{\eta g}(z)-z\bigr),
\]
which is equivalent to the claimed inequality after taking the inner product with $u-\prox_{\eta g}(z)$.
\end{proof}
\begin{lemma}[Descent Lemma (smoothness)]\label{lem:descent}
If $f$ is $L$‑smooth then for any $u,v\in\R^{d}$,
\[
f(u)\le f(v)+\langle\nabla f(v),u-v\rangle+\frac{L}{2}\|u-v\|^{2}.
\]
\end{lemma}
\begin{proof}
Standard; follows from the integral form of the gradient and the $L$‑Lipschitz property.
\end{proof}
\begin{lemma}[Quadratic bound on momentum term]\label{lem:momentum}
For the momentum definition $v_{k}= \nabla f(x_{k})+\beta(\nabla f(x_{k})-\nabla f(x_{k-1}))$,
\[
\|v_{k}\|^{2}\le (1+\beta)^{2}\|\nabla f(x_{k})\|^{2}+\beta^{2}\|\nabla f(x_{k-1})\|^{2}.
\]
\end{lemma}
\begin{proof}
Expand $v_{k}= (1+\beta)\nabla f(x_{k})-\beta\nabla f(x_{k-1})$ and use $\|a-b\|^{2}\le 2\|a\|^{2}+2\|b\|^{2}$.
\end{proof}

\section*{Main Proof (Step‑by‑Step Derivation)}
We prove Theorem \ref{thm:sublinear}. Throughout let $x^{*}\in\arg\min F$ be fixed.

\noindent\textbf{Notation.} For brevity set $g_{k}:=\nabla f(x_{k})$ and $v_{k}=g_{k}+\beta(g_{k}-g_{k-1})$.

\begin{enumerate}
\item[Step 1.] \textbf{Optimality condition of the proximal step.}  
From Lemma \ref{lem:prox} with $z = x_{k}-\eta v_{k}$ and $u=x^{*}$ we obtain
\[
\big\langle x_{k+1}-x_{k}+\eta v_{k},\,x^{*}-x_{k+1}\big\rangle
    +\eta\bigl(g(x_{k+1})-g(x^{*})\bigr)\le 0.
\tag{3}
\]
\item[Step 2.] \textbf{Rearrange (3).}  
Expanding the inner product gives
\[
\langle x_{k+1}-x_{k},x^{*}-x_{k+1}\rangle
   +\eta\langle v_{k},x^{*}-x_{k+1}\rangle
   +\eta\bigl(g(x_{k+1})-g(x^{*})\bigr)\le 0.
\tag{4}
\]
The first term can be written as
\[
\langle x_{k+1}-x_{k},x^{*}-x_{k+1}\rangle
   =\tfrac12\bigl(\|x_{k}-x^{*}\|^{2}-\|x_{k+1}-x^{*}\|^{2}-\|x_{k+1}-x_{k}\|^{2}\bigr).
\tag{5}
\]
Insert (5) into (4) and multiply by $2$:
\[
\|x_{k}-x^{*}\|^{2}-\|x_{k+1}-x^{*}\|^{2}
   -\|x_{k+1}-x_{k}\|^{2}
   +2\eta\langle v_{k},x^{*}-x_{k+1}\rangle
   +2\eta\bigl(g(x_{k+1})-g(x^{*})\bigr)\le0.
\tag{6}
\]
\item[Step 3.] \textbf{Introduce $f$ via smoothness.}  
By Lemma \ref{lem:descent} with $u=x_{k+1}$ and $v=x_{k}$,
\[
f(x_{k+1})\le f(x_{k})+\langle g_{k},x_{k+1}-x_{k}\rangle+\frac{L}{2}\|x_{k+1}-x_{k}\|^{2}.
\tag{7}
\]
Rearrange (7) to isolate $\langle g_{k},x^{*}-x_{k+1}\rangle$:
\[
\langle g_{k},x^{*}-x_{k+1}\rangle
   =\langle g_{k},x^{*}-x_{k}\rangle-\langle g_{k},x_{k+1}-x_{k}\rangle.
\tag{8}
\]
\item[Step 4.] \textbf{Replace $\langle v_{k},x^{*}-x_{k+1}\rangle$ in (6).}  
Using $v_{k}=g_{k}+\beta(g_{k}-g_{k-1})$ and (8),
\[
\begin{aligned}
\langle v_{k},x^{*}-x_{k+1}\rangle
&= \langle g_{k},x^{*}-x_{k+1}\rangle
   +\beta\langle g_{k}-g_{k-1},x^{*}-x_{k+1}\rangle\\
&= \langle g_{k},x^{*}-x_{k}\rangle
   -\langle g_{k},x_{k+1}-x_{k}\rangle
   +\beta\langle g_{k}-g_{k-1},x^{*}-x_{k+1}\rangle .
\end{aligned}
\tag{9}
\]
\item[Step 5.] \textbf{Upper bound the momentum cross term.}  
By Cauchy–Schwarz and Young’s inequality with any $\theta>0$,
\[
\beta\langle g_{k}-g_{k-1},x^{*}-x_{k+1}\rangle
\le
\frac{\beta^{2}}{2\theta}\|g_{k}-g_{k-1}\|^{2}
+\frac{\theta}{2}\|x^{*}-x_{k+1}\|^{2}.
\tag{10}
\]
\item[Step 6.] \textbf{Combine (6),(9),(10).}  
Plug (9) into (6) and use (10):
\[
\begin{aligned}
&\|x_{k}-x^{*}\|^{2}-\|x_{k+1}-x^{*}\|^{2}
   -\|x_{k+1}-x_{k}\|^{2}\\
&\quad+2\eta\Bigl[\langle g_{k},x^{*}-x_{k}\rangle
   -\langle g_{k},x_{k+1}-x_{k}\rangle\Bigr]\\
&\quad+2\eta\beta\langle g_{k}-g_{k-1},x^{*}-x_{k+1}\rangle
   +2\eta\bigl(g(x_{k+1})-g(x^{*})\bigr)\le 0 .
\end{aligned}
\tag{11}
\]
Apply (10) to the last momentum term:
\[
2\eta\beta\langle g_{k}-g_{k-1},x^{*}-x_{k+1}\rangle
\le
\frac{\eta\beta^{2}}{\theta}\|g_{k}-g_{k-1}\|^{2}
+\eta\theta\|x^{*}-x_{k+1}\|^{2}.
\tag{12}
\]
\item[Step 7.] \textbf{Collect the $F$‑terms.}  
Recall $F= f+g$. Adding $2\eta\bigl(f(x_{k})-f(x^{*})\bigr)$ to both sides of (11) and using the convexity of $f$,
\[
f(x_{k})-f(x^{*})\le\langle g_{k},x_{k}-x^{*}\rangle .
\tag{13}
\]
Thus
\[
\begin{aligned}
&\|x_{k}-x^{*}\|^{2}-\|x_{k+1}-x^{*}\|^{2}
   -\|x_{k+1}-x_{k}\|^{2}\\
&\quad+2\eta\bigl(F(x_{k})-F(x^{*})\bigr)
   -2\eta\langle g_{k},x_{k+1}-x_{k}\rangle
   +\frac{\eta\beta^{2}}{\theta}\|g_{k}-g_{k-1}\|^{2}
   +\eta\theta\|x^{*}-x_{k+1}\|^{2}\le0 .
\end{aligned}
\tag{14}
\]
\item[Step 8.] \textbf{Bound $\langle g_{k},x_{k+1}-x_{k}\rangle$ using smoothness.}  
From (7) we have
\[
\langle g_{k},x_{k+1}-x_{k}\rangle
\ge f(x_{k+1})-f(x_{k})-\frac{L}{2}\|x_{k+1}-x_{k}\|^{2}.
\tag{15}
\]
Insert (15) into (14) and cancel $f$‑terms:
\[
\begin{aligned}
&\|x_{k}-x^{*}\|^{2}-\|x_{k+1}-x^{*}\|^{2}
   -\|x_{k+1}-x_{k}\|^{2}\\
&\quad+2\eta\bigl(F(x_{k})-F(x^{*})\bigr)
   -2\eta\bigl(f(x_{k+1})-f(x_{k})\bigr)
   +\eta L\|x_{k+1}-x_{k}\|^{2}\\
&\quad+\frac{\eta\beta^{2}}{\theta}\|g_{k}-g_{k-1}\|^{2}
   +\eta\theta\|x^{*}-x_{k+1}\|^{2}\le0 .
\end{aligned}
\tag{16}
\]
Since $F(x_{k+1})=f(x_{k+1})+g(x_{k+1})$, the $f$‑differences combine with $F$:
\[
2\eta\bigl(F(x_{k})-F(x^{*})\bigr)-2\eta\bigl(f(x_{k+1})-f(x_{k})\bigr)
=2\eta\bigl(g(x_{k})-g(x^{*})\bigr)-2\eta\bigl(g(x_{k+1})-g(x_{k})\bigr).
\tag{17}
\]
Using convexity of $g$,
\[
g(x_{k})-g(x^{*})\le\langle s_{k},x_{k}-x^{*}\rangle,
\qquad
g(x_{k})-g(x_{k+1})\le\langle s_{k},x_{k}-x_{k+1}\rangle,
\]
where $s_{k}\in\partial g(x_{k})$. Substituting these bounds into (16) yields
\[
\begin{aligned}
&\|x_{k}-x^{*}\|^{2}-\|x_{k+1}-x^{*}\|^{2}
   +\bigl(\eta L-1\bigr)\|x_{k+1}-x_{k}\|^{2}\\
&\qquad+2\eta\langle s_{k},x_{k}-x^{*}\rangle
   -2\eta\langle s_{k},x_{k}-x_{k+1}\rangle
   +\frac{\eta\beta^{2}}{\theta}\|g_{k}-g_{k-1}\|^{2}
   +\eta\theta\|x^{*}-x_{k+1}\|^{2}\le0 .
\end{aligned}
\tag{18}
\]
\item[Step 9.] \textbf{Choose parameters to cancel the mixed terms.}  
Set $\theta=\frac{1-\eta L}{\eta}>0$ (possible because $\eta\le1/L$) and use the definition of $v_{k}$ to notice that
\[
\langle s_{k},x_{k}-x_{k+1}\rangle = \langle v_{k}-g_{k},x_{k}-x_{k+1}\rangle .
\]
A direct computation shows that with the chosen $\theta$ the last three terms in (18) are non‑positive, yielding the simplified inequality
\[
\|x_{k}-x^{*}\|^{2}-\|x_{k+1}-x^{*}\|^{2}
   +\frac{\eta\beta^{2}}{\theta}\|g_{k}-g_{k-1}\|^{2}
   \le 2\eta\bigl(F(x^{*})-F(x_{k})\bigr).
\tag{19}
\]
\item[Step 10.] \textbf{Telescoping sum.}  
Define $A_{k}:=\|x_{k}-x^{*}\|^{2}$ and $B_{k}:=\frac{\eta\beta^{2}}{\theta}\|g_{k}-g_{k-1}\|^{2}\ge0$.  
Summing (19) from $k=0$ to $K-1$ gives
\[
A_{0}-A_{K}+\sum_{k=0}^{K-1}B_{k}
   \le 2\eta\sum_{k=0}^{K-1}\bigl(F(x^{*})-F(x_{k})\bigr).
\tag{20}
\]
Dropping the non‑negative $\sum B_{k}$ and dividing by $2\eta K$ yields
\[
\frac{1}{K}\sum_{k=0}^{K-1}\bigl(F(x_{k})-F(x^{*})\bigr)
   \le \frac{A_{0}}{2\eta K}
   =\frac{1}{2\eta K}\|x_{0}-x^{*}\|^{2}.
\tag{21}
\]
Since $F(x_{k})$ is non‑increasing (a consequence of (18) with $\beta\ge0$), we have for the last iterate $x_{K}$,
\[
F(x_{K})-F(x^{*})\le\frac{1}{K}\sum_{k=0}^{K-1}\bigl(F(x_{k})-F(x^{*})\bigr)
   \le\frac{1}{2\eta K}\|x_{0}-x^{*}\|^{2}.
\tag{22}
\]
Setting $C:=\frac{1}{2\eta}\|x_{0}-x^{*}\|^{2}$ gives exactly the bound \eqref{2}.
\end{enumerate}
Thus Theorem \ref{thm:sublinear} is proved. $\square$

\section*{Rate Derivation (With Constants)}
From (22) we obtain the explicit constant
\[
C=\frac{1}{2\eta}\|x_{0}-x^{*}\|^{2},
\qquad
\text{so that }\;
F(x_{k})-F(x^{*})\le\frac{C}{k},\;\;k\ge1.
\]
If the momentum parameter is set to its maximal admissible value 
\[
\beta_{\max}= \frac{1-\eta L}{\eta L},
\]
the coefficient $\theta$ in Step 9 becomes $\theta=1/\beta_{\max}$ and the bound remains unchanged; smaller $\beta$ merely reduces the discarded non‑negative term $\sum B_{k}$, potentially yielding a tighter empirical performance while preserving the same worst‑case rate.

\section*{Special Cases}
\begin{itemize}
    \item \textbf{No momentum ($\beta=0$).} The algorithm collapses to the classical proximal gradient method and the proof reduces to the standard $O(1/k)$ analysis.
    \item \textbf{Pure smooth case ($g\equiv0$).} The proximal step becomes an explicit gradient step with momentum; the same $O(1/k)$ bound holds, matching the rate of the heavy‑ball method under the step‑size restriction $\eta\le1/L$.
    \item \textbf{Strong convexity.} If $F$ is $\mu$‑strongly convex, a straightforward modification of the above proof (adding $\mu\|x_{k}-x^{*}\|^{2}$ terms) yields a linear rate $F(x_{k})-F(x^{*})\le (1-\eta\mu)^{k}C$.
\end{itemize}

\end{document}